\documentclass[10pt]{article}
\usepackage[a4paper, margin=0.75in]{geometry}
\usepackage{titlesec}
\usepackage{booktabs}
\usepackage{tabularx}
\usepackage{xcolor}
\usepackage{enumitem}
\usepackage{tgheros}
\usepackage[T1]{fontenc}
\usepackage{array}
\usepackage{textcomp}
\usepackage[hidelinks]{hyperref}
\usepackage{tikz}
\usetikzlibrary{shapes.rounded, arrows.meta, positioning, calc, fit}
\renewcommand{\familydefault}{\sfdefault}

%% ===== Colour palette =====
\definecolor{primary}{RGB}{15,52,96}
\definecolor{accent}{RGB}{198,124,0}
\definecolor{tier1}{RGB}{20,100,60}
\definecolor{tier2}{RGB}{140,80,0}
\definecolor{tier3}{RGB}{80,80,80}
\definecolor{lightgrey}{gray}{0.88}
\definecolor{midgrey}{gray}{0.65}
\definecolor{darkgrey}{gray}{0.38}

%% ===== Section style =====
\titleformat{\section}{%
  \color{primary}\bfseries\large\raggedright
}{}{0em}{}[\color{lightgrey}{\titlerule[1.5pt]}\vspace{-4pt}]

\titleformat{\subsection}{%
  \color{accent}\bfseries\normalsize\raggedright
}{}{0em}{}

\titlespacing*{\section}{0pt}{18pt}{8pt}
\titlespacing*{\subsection}{0pt}{12pt}{4pt}

\setlength{\parindent}{0pt}
\setlength{\parskip}{6pt}
\renewcommand{\arraystretch}{1.25}
\hypersetup{colorlinks=true, urlcolor=primary}

\newcommand{\spacer}{\vspace{10pt}}
\newcommand{\bigspacer}{\vspace{16pt}}

%% =====================================================
\begin{document}

%% ----- Title block -----
\begin{center}
  \vspace*{0.1cm}
  {\color{primary}\LARGE\bfseries NEWCROSS AND PAN OCEAN GROUP}\\[6pt]
  {\color{primary}\large\bfseries AI and Data Roadmap --- Practical Intelligence
   for Nigerian Oil and Gas Operations}\\[10pt]
  {\color{darkgrey}\normalsize Version 2.0\quad\textbar\quad May 2026\quad\textbar\quad
   Companion to the Strategic Vision Paper}
\end{center}

\bigspacer

\noindent\colorbox{primary!8}{%
\begin{minipage}{\dimexpr\textwidth-2\fboxsep}
  \vspace{6pt}
  \textbf{\color{primary}Guiding Principle.}\quad
  Don't build AI for AI's sake. Build AI for Nigerian realities. Every initiative
  in this document passes three tests before it is considered viable:

  \vspace{4pt}
  \textbf{(1)} Does it solve a real, costly problem at the group today? \\
  \textbf{(2)} Can it be built with data and infrastructure that actually exists? \\
  \textbf{(3)} Could another Nigerian operator pay for it within three years?

  \vspace{4pt}
  If any answer is no, the initiative is deprioritised.
  \vspace{6pt}
\end{minipage}}

\bigspacer

\noindent\colorbox{accent!8}{%
\begin{minipage}{\dimexpr\textwidth-2\fboxsep}
  \vspace{6pt}
  \textbf{\color{accent}Note on Sequencing.}\quad
  The tiers in this document represent \emph{data and infrastructure readiness},
  not a calendar plan. Tier 1 initiatives can begin with the data and instrumentation
  that realistically exists today. Tier 2 requires the data foundation built during
  Tier 1 to be in place. Tier 3 requires sustained investment in data quality and
  governance over twelve to eighteen months before machine learning is viable.

  \vspace{4pt}
  Specific dates are deliberately not assigned. The pace of execution depends on
  internal commitment, parallel initiatives the external view cannot fully see,
  and the group's own prioritisation.
  \vspace{6pt}
\end{minipage}}

\bigspacer

%% ----------------------------------------------------------
\section{Nigerian Operational Reality --- Design Constraints}
%% ----------------------------------------------------------

Every AI system built for this group must be designed around these realities
from day one. They are not edge cases. They are the standard operating environment.

\spacer
\small
\begin{tabularx}{\columnwidth}{@{}lX@{}}
\toprule
\textbf{Operational Reality} & \textbf{Design Implication} \\
\midrule
Intermittent internet at remote sites (Ogharefe, Owa Aladinma, AEPP corridor)
  & Offline-first; edge inference where possible; sync when connected \\
Power instability across operational sites
  & Async processing; graceful degradation; no cloud-only dependency \\
Low bandwidth at field locations
  & Lightweight models; local inference; minimal data transfer \\
NUPRC regulatory framework including April 2026 GHG directive
  & Compliance logic baked into data models, not added as an afterthought \\
PIA 2021 HCDT obligations (3\% of opex per entity)
  & HCDT tracking integrated into procurement and finance data pipelines
    from the first day of data collection \\
Nigerian Content Act requirements
  & Spend and employment data structured for NCA reporting from day one \\
Naira currency and FX exposure
  & Price and operate in Naira; minimise dollar-denominated cloud costs \\
Data quality legacy issues across SAP and operational systems
  & Assume dirty data; build cleaning and validation pipelines before ML \\
Trust and explainability culture in operations teams
  & AI recommendations must explain their reasoning to be accepted by
    engineers and plant managers \\
Multi-entity group structure (three legal entities, partly shared operations)
  & All systems must support entity-level and consolidated views simultaneously \\
\bottomrule
\end{tabularx}
\normalsize

\bigspacer

%% ----------------------------------------------------------
\section{Tier 1 --- Achievable with Existing Data and Infrastructure}
%% ----------------------------------------------------------

\noindent\colorbox{tier1!8}{%
\begin{minipage}{\dimexpr\textwidth-2\fboxsep}
  \vspace{4pt}
  {\small\color{tier1}\textbf{Tier 1 Criteria:} Real problem with immediate cost
  impact. Buildable with data and infrastructure that exists today, plus
  realistic instrumentation. Commercial potential clear.}
  \vspace{4pt}
\end{minipage}}

\bigspacer

\subsection*{1. Predictive Maintenance --- OOGP Gas Plant (Highest Priority)}

\textbf{Problem.}
The Ovade-Ogharefe Gas Plant is the group's most complex and highest-consequence
asset. It is a large-scale cryogenic gas processing facility with 29 LPG storage
tanks. Equipment failures --- compressors, cryogenic units, storage pressure
systems --- are extremely costly and potentially dangerous. Maintenance is
currently reactive: the problem is discovered after it has already become a
failure or near-failure.

\textbf{Solution.}
Instrument compressors, cryogenic processing units, and LPG storage tanks
with vibration, temperature, and pressure sensors. Train time-series anomaly
detection models on sensor readings combined with existing SAP maintenance
logs. Output: failure probability scores per asset class, recommended
maintenance windows, escalation alerts to the right personnel.

\textbf{Why now.}
The Phase 2 cryogenic expansion was completed in 2022. Maintenance pattern
data from the new cryogenic systems is still building. This is the ideal
moment to instrument fully and establish clean data before failure patterns
are entrenched.

\spacer
\small
\begin{tabularx}{\columnwidth}{@{}lX@{}}
\toprule
\textbf{Item} & \textbf{Detail} \\
\midrule
Data required     & SAP maintenance logs; SCADA feeds (existing); equipment OEM
                    specs and fault codes; sensor readings (new instrumentation
                    required) \\
Tech stack        & Python, scikit-learn / PyTorch, time-series libraries,
                    InfluxDB or TimescaleDB for sensor data, Azure IoT Hub for
                    ingestion \\
Expected ROI      & Reduction in unplanned downtime (one avoided plant shutdown
                    justifies the full project); extended equipment lifespan;
                    reduction in emergency parts procurement \\
Commercial product & \textbf{FieldMaintain} SaaS --- every gas plant operator in
                    Nigeria has this problem. OOGP is the reference deployment. \\
\bottomrule
\end{tabularx}
\normalsize

\bigspacer

\subsection*{2. Regulatory Reporting Automation --- NUPRC Returns and GHG Directive}

\textbf{Problem.}
Quarterly and annual NUPRC regulatory submissions require pulling data from
multiple systems (SAP, spreadsheets, manual records), formatting it to NUPRC
templates, and submitting under tight deadlines. This takes teams days per
submission cycle across OML-98, OML-147, and OOGP. The April 2026 NUPRC
directive adds mandatory measurement-based methane and GHG reporting templates,
making manual compliance even more operationally demanding.

\textbf{Solution.}
Build data extraction pipelines from SAP and operational systems. Use
structured template-filling with LLM assistance to auto-populate required
report sections, including the new GHG/methane templates. Add validation rules
against NUPRC submission requirements. Output: draft submissions ready for
human review and sign-off. For GHG reporting, sensor data feeds directly into
the reporting pipeline.

\spacer
\small
\begin{tabularx}{\columnwidth}{@{}lX@{}}
\toprule
\textbf{Item} & \textbf{Detail} \\
\midrule
Data required     & SAP data; existing NUPRC templates; past approved
                    submissions (for validation); OOGP sensor data for GHG
                    measurement \\
Tech stack        & Python, LangChain or direct LLM API (locally hosted for
                    sensitive data), SAP data connectors, PDF generation
                    library \\
Expected ROI      & Days $\to$ hours per submission cycle; fewer errors;
                    full audit trail; GHG directive compliance without
                    additional headcount \\
Commercial product & \textbf{RegulatoryReport} --- painful for every operator;
                    foreign vendors consistently fail on Nigerian regulatory
                    specifics \\
\bottomrule
\end{tabularx}
\normalsize

\bigspacer

\subsection*{3. Nigerian Content and HCDT Compliance Tracker}

\textbf{Problem.}
Tracking, calculating, and reporting Nigerian Content Act compliance (local
spend percentage, local employment, training documentation) is manual,
error-prone, and creates significant audit risk. Under PIA 2021, the HCDT
obligation (3\% of opex) adds a further layer of reporting across all three
group entities --- each with distinct cost bases and procurement records.

\textbf{Solution.}
Integrate with SAP procurement and HR data across all three entities.
Auto-classify suppliers and spend by Nigerian Content category using NLP
classification models. Calculate HCDT contributions per entity from opex
data. Generate compliance reports in NCA-required format with early-warning
flags before audit cycles.

\spacer
\small
\begin{tabularx}{\columnwidth}{@{}lX@{}}
\toprule
\textbf{Item} & \textbf{Detail} \\
\midrule
Data required     & SAP procurement records; HR employee and training data;
                    vendor master lists; opex ledger per entity \\
Tech stack        & Python, spaCy or fine-tuned BERT for text classification,
                    SAP data connectors, report generation pipeline \\
Expected ROI      & Audit readiness; reduced compliance risk; staff time
                    savings; HCDT obligations tracked automatically \\
Commercial product & \textbf{NigerianContent.ai} --- mandatory for all operators;
                    PIA 2021 raised the stakes significantly \\
\bottomrule
\end{tabularx}
\normalsize

\bigspacer

\subsection*{4. AEPP Integrity and Third-Party Tariff Reconciliation}

\textbf{Problem.}
The Amukpe-Escravos Pipeline serves third-party injectors on a tariff basis,
meaning metering accuracy and theft detection are direct revenue issues, not
merely operational ones. The 67 km corridor through the Niger Delta with CHDD
construction still presents exposure. Tariff reconciliation for third-party
injectors is currently paper-intensive and manually reconciled.

\textbf{Solution.}
Deploy flow meter anomaly detection across all injection and delivery points.
Satellite imagery change detection (monthly comparison) along the 67 km Right
of Way to identify new construction, dig activity, or vegetation clearance
near the pipeline. Automated metering and tariff billing for injector volumes
with full audit trail. GPS-coordinated alert system for field response teams.

\spacer
\small
\begin{tabularx}{\columnwidth}{@{}lX@{}}
\toprule
\textbf{Item} & \textbf{Detail} \\
\midrule
Data required     & SCADA flow meter readings; AEPP GPS coordinates;
                    satellite imagery (NASRDA preferred as Nigerian-first;
                    Planet Labs or Airbus Defence as alternatives); existing
                    tariff and injection records from SAP \\
Tech stack        & Python, GDAL/rasterio for satellite imagery, statistical
                    process control for flow anomaly detection, GIS integration \\
Expected ROI      & Crude theft prevention; metering accuracy directly protects
                    tariff revenue; removes manual reconciliation overhead \\
Commercial product & \textbf{FieldWatch} --- every operator with long-haul Niger
                    Delta pipelines faces identical problems; AEPP is a 67 km
                    reference site \\
\bottomrule
\end{tabularx}
\normalsize

\bigspacer

\subsection*{5. HSE Digital Permit-to-Work and Incident Reporting}

\textbf{Problem.}
OOGP operates cryogenic processing equipment and 29 LPG storage tanks in a
high-consequence environment. Permit-to-work (PTW) systems are paper-based.
Incident reporting and near-miss tracking are manual. NUPRC has increasing
expectations around HSE reporting standards for high-consequence facilities.

\textbf{Solution.}
Replace paper PTW with a digital system configured for the OOGP environment.
Integrate incident reporting and near-miss tracking with a searchable
database. Generate HSE reports automatically for NUPRC submissions. Mobile-
accessible for field personnel; offline-capable for areas without reliable
connectivity.

\spacer
\small
\begin{tabularx}{\columnwidth}{@{}lX@{}}
\toprule
\textbf{Item} & \textbf{Detail} \\
\midrule
Data required     & Existing PTW records (digitise backlog for training);
                    incident and near-miss historical records; NUPRC HSE
                    reporting templates \\
Tech stack        & React Native (mobile + offline-first), FastAPI backend,
                    PostgreSQL, Azure sync \\
Expected ROI      & Liability reduction; NUPRC compliance readiness; searchable
                    incident history for trend analysis \\
Commercial product & \textbf{HSETrack} --- foreign HSE software does not
                    understand Nigerian site conditions or NUPRC requirements \\
\bottomrule
\end{tabularx}
\normalsize

\bigspacer

%% ----------------------------------------------------------
\section{Tier 2 --- Build After Tier 1 Data Foundation Is in Place}
%% ----------------------------------------------------------

\noindent\colorbox{tier2!8}{%
\begin{minipage}{\dimexpr\textwidth-2\fboxsep}
  \vspace{4pt}
  {\small\color{tier2}\textbf{Tier 2 Criteria:}
  High strategic value but requires data pipeline work from Tier 1 to be in
  place first. Begin planning in parallel; begin building once data foundations
  are established.}
  \vspace{4pt}
\end{minipage}}

\bigspacer

\subsection*{6. Document Intelligence --- Well Logs, Drilling Reports, Contracts}

\textbf{Problem.}
OML-98 has been producing since 1976. That is fifty years of drilling reports,
well logs, geological surveys, production records, and contracts --- most of
which exist as PDFs, scanned paper, or disconnected digital files. Finding
specific data (a well's pressure history, a vendor's contract terms, a drilling
decision rationale) takes hours or days. The group's institutional memory is
largely inaccessible.

\textbf{Solution.}
OCR pipeline to digitise historical paper documents. LLM-based document
classification and metadata extraction. Semantic search (RAG architecture)
so engineers can query natural-language questions against the full document
archive. Unlock fifty years of operational knowledge.

\spacer
\small
\begin{tabularx}{\columnwidth}{@{}lX@{}}
\toprule
\textbf{Item} & \textbf{Detail} \\
\midrule
Data required     & Existing document archives (digital and physical).
                    Digitisation of paper records is the primary prerequisite. \\
Tech stack        & Python, Tesseract / AWS Textract for OCR, LangChain
                    + vector database (pgvector or Weaviate), locally-hosted
                    LLM (Llama3 or Mistral) for sensitive data \\
Security note     & Well data and production history is commercially sensitive.
                    Use locally-hosted LLMs only for this application. \\
Expected ROI      & Engineer productivity; faster decision-making;
                    institutional knowledge preserved beyond individual tenure \\
\bottomrule
\end{tabularx}
\normalsize

\bigspacer

\subsection*{7. Unified Operations Dashboard with Anomaly Detection}

\textbf{Problem.}
Management has no real-time view across OML-98, OML-147, OOGP, and AEPP
operations simultaneously. Reports are generated periodically, so problems
surface late. The three-entity group structure means consolidated visibility
is even harder to achieve manually.

\textbf{Solution.}
Unified data integration layer pulling from SAP, SCADA, and operational
systems across all entities. Real-time dashboard (web and mobile-accessible)
with configurable KPIs per entity and consolidated group view. Anomaly
detection layer flagging deviations from expected parameters. Alert system
with escalation logic.

\spacer
\small
\begin{tabularx}{\columnwidth}{@{}lX@{}}
\toprule
\textbf{Item} & \textbf{Detail} \\
\midrule
Data required     & SCADA feeds; SAP operational data; logistics and
                    dispatch records; gas plant throughput data; AEPP
                    metering records \\
Tech stack        & Python (FastAPI backend), React or React Native (frontend),
                    TimescaleDB, statistical anomaly detection layer \\
Expected ROI      & Faster management response to operational issues;
                    group-level visibility that currently does not exist
                    in any form \\
Commercial product & \textbf{WellWatch} SaaS --- configurability for
                    multi-entity, multi-block operations is the key
                    differentiator \\
\bottomrule
\end{tabularx}
\normalsize

\bigspacer

\subsection*{8. Oil Marketing Intelligence}

\textbf{Problem.}
Crude lifting schedules, buyer allocation, and NNPC PSC equity crude
reconciliation on OML-147 are managed manually. The domestic crude market
has materially changed with Dangote refinery's entry. Oil marketing teams
are making significant commercial decisions without a systematic data
environment to support them.

\textbf{Solution.}
Lifting schedule management and buyer allocation tracking system. PSC equity
crude reconciliation with NNPC, integrating with the OML-147 PSC reporting
pipeline. Domestic and export pricing dashboard tracking refinery netbacks,
regional differentials, and historical buyer performance.

\spacer
\small
\begin{tabularx}{\columnwidth}{@{}lX@{}}
\toprule
\textbf{Item} & \textbf{Detail} \\
\midrule
Data required     & Lifting records; buyer contracts; NNPC PSC equity
                    crude schedules; Platts or Argus price feeds (API) \\
Tech stack        & Python backend, React frontend, PostgreSQL, price
                    feed integrations \\
Expected ROI      & Lifting schedule accuracy; reduced reconciliation
                    disputes; better-informed commercial decisions in
                    a shifting market \\
\bottomrule
\end{tabularx}
\normalsize

\bigspacer

%% ----------------------------------------------------------
\section{Tier 3 --- Requires Sustained Data Foundation}
%% ----------------------------------------------------------

\noindent\colorbox{tier3!10}{%
\begin{minipage}{\dimexpr\textwidth-2\fboxsep}
  \vspace{4pt}
  {\small\color{tier3}\textbf{Tier 3 Criteria:}
  High long-term value but requires sustained upstream data pipeline work
  before machine learning is viable. Plan architecture now; do not attempt
  to build now.}
  \vspace{4pt}
\end{minipage}}

\bigspacer

\subsection*{9. Production Optimisation --- OML-98 and OML-147}

\textbf{Problem.}
Well and field performance across both OML blocks is managed based on
engineering experience and periodic review rather than continuous
optimisation. Small improvements in choke settings, injection schedules,
and pressure management across two producing blocks translate to meaningful
revenue.

\textbf{Solution.}
Optimisation models trained on historical production data. Recommendations
for well scheduling, injection rates, and choke settings. Simulation
environment for testing operational decisions without field risk.

\textbf{Prerequisite.}
Requires significant data quality improvement work from Tier 1 and Tier 2
initiatives. This is a sustained build, not a starting point.

\bigspacer

\subsection*{10. Procurement and Supply Chain Optimisation}

\textbf{Problem.}
Procurement for remote field operations involves significant lead times,
FX exposure on imported parts, and frequent emergency orders at premium cost.
Vendor performance is not systematically tracked. Three entities with partly
overlapping procurement needs creates further inefficiency.

\textbf{Solution.}
Demand forecasting for spare parts based on maintenance schedules and
equipment age (feeding directly from the Tier 1 predictive maintenance
system). Vendor risk scoring from historical delivery performance. FX
exposure dashboard for imported goods. Consolidated group procurement view
to identify cross-entity opportunities.

\bigspacer

%% ----------------------------------------------------------
\section{AI Design Principles for Group Deployments}
%% ----------------------------------------------------------

\small
\begin{tabularx}{\columnwidth}{@{}lX@{}}
\toprule
\textbf{Principle} & \textbf{Rationale and Application} \\
\midrule
Offline-first
  & Core AI functions must produce useful output without reliable internet.
    OOGP, Ogharefe, and the AEPP corridor are not reliably connected
    environments. \\
Edge inference where possible
  & Run model inference locally at the asset; send only anomaly alerts and
    summaries to the cloud. Reduces bandwidth, latency, and cloud cost. \\
Explainable outputs
  & Every AI recommendation must include a plain-language explanation.
    ``Compressor Unit 3: 78\% failure probability within 14 days based on
    bearing vibration trend and last maintenance date.'' Not just a score. \\
Locally hosted for sensitive data
  & Well data, production history, commercial contracts, and lifting records
    must not leave the group's Azure tenant or on-premise servers. Use
    locally-hosted LLMs (Llama3, Mistral) for these applications. \\
Assume dirty data
  & Every pipeline includes a validation and cleaning layer before data
    reaches any model. Data quality is logged and reported, not silently
    fixed. \\
Multi-entity from day one
  & Every model and dashboard is built to operate at both entity and
    consolidated group level. Single-entity models will never be commercially
    reusable. \\
Compliance logic is first-class
  & NUPRC reporting requirements, HCDT calculations, and NCA classifications
    are embedded in data models, not bolted on as reporting features later. \\
\bottomrule
\end{tabularx}
\normalsize

\bigspacer

%% ----------------------------------------------------------
\section{Sequencing Logic --- Dependencies, Not Dates}
%% ----------------------------------------------------------

The diagram below shows how the ten initiatives depend on one another. It
is read from left to right --- arrows mean ``cannot start meaningfully until.''
The pace is set by internal commitment, not by this document.

\spacer

\begin{center}
\begin{tikzpicture}[
  node distance=0.4cm and 0.7cm,
  found/.style={draw=primary, fill=primary!8, rounded corners=4pt,
                minimum width=3.6cm, minimum height=1.0cm,
                align=center, font=\small\bfseries, text=black,
                inner sep=4pt},
  t1/.style={draw=tier1!70, fill=tier1!10, rounded corners=4pt,
             minimum width=3.6cm, minimum height=1.0cm,
             align=center, font=\small, inner sep=4pt},
  t2/.style={draw=tier2!70, fill=tier2!10, rounded corners=4pt,
             minimum width=3.6cm, minimum height=1.0cm,
             align=center, font=\small, inner sep=4pt},
  t3/.style={draw=tier3!70, fill=tier3!10, rounded corners=4pt,
             minimum width=3.6cm, minimum height=1.0cm,
             align=center, font=\small, inner sep=4pt},
  arr/.style={-{Stealth[length=4pt]}, thick, color=midgrey},
]

%% Foundations column
\node[found] (audit)  {Connectivity \&\\Data Audit};
\node[found, below=0.5cm of audit] (gov) {Data Governance\\Framework};
\node[found, below=0.5cm of gov] (instr) {OOGP Sensor\\Instrumentation};
\node[found, below=0.5cm of instr] (sap) {SAP Data\\Connectors};

%% Tier 1 column
\node[t1, right=2.0cm of audit] (i2) {2. Regulatory\\Reporting + GHG};
\node[t1, below=0.5cm of i2] (i3) {3. NCA / HCDT\\Tracker};
\node[t1, below=0.5cm of i3] (i1) {1. Predictive\\Maintenance OOGP};
\node[t1, below=0.5cm of i1] (i4) {4. AEPP Integrity\\\& Tariff};
\node[t1, below=0.5cm of i4] (i5) {5. HSE Digital\\PTW};

%% Tier 2 column
\node[t2, right=2.0cm of i2] (i7) {7. Unified Ops\\Dashboard};
\node[t2, below=0.5cm of i7] (i8) {8. Oil Marketing\\Intelligence};
\node[t2, below=0.5cm of i8] (i6) {6. Document\\Intelligence};

%% Tier 3 column
\node[t3, right=2.0cm of i7] (i9) {9. Production\\Optimisation};
\node[t3, below=0.5cm of i9] (i10) {10. Procurement\\\& Supply Chain};

%% Arrows: foundations -> Tier 1
\draw[arr] (audit.east) -- (i2.west);
\draw[arr] (audit.east) -- (i3.west);
\draw[arr] (instr.east) -- (i1.west);
\draw[arr] (sap.east) -- (i3.west);
\draw[arr] (sap.east) -- (i4.west);

%% Arrows: Tier 1 -> Tier 2
\draw[arr] (i1.east) -- (i7.west);
\draw[arr] (i4.east) -- (i7.west);
\draw[arr] (i2.east) -- (i8.west);

%% Arrows: Tier 2 -> Tier 3
\draw[arr] (i7.east) -- (i9.west);
\draw[arr] (i1.east) -- ([yshift=-2pt] i10.west);

%% Headers
\node[font=\small\bfseries, color=primary, above=0.2cm of audit] {Foundations};
\node[font=\small\bfseries, color=tier1, above=0.2cm of i2] {Tier 1};
\node[font=\small\bfseries, color=tier2, above=0.2cm of i7] {Tier 2};
\node[font=\small\bfseries, color=tier3, above=0.2cm of i9] {Tier 3};

\end{tikzpicture}
\end{center}

\spacer
{\small\color{darkgrey}
\textbf{Reading the diagram.}\enspace
Foundations are the prerequisites that make everything else possible. Tier 1
initiatives can begin as soon as their specific foundations are in place ---
not all at once. Tier 2 follows when Tier 1 has produced clean, reliable data.
Tier 3 follows when Tier 2 dashboards and document intelligence have matured
enough to support model training. There is no fixed calendar; the sequence is
the constraint, not the dates.
}

\bigspacer

%% ----------------------------------------------------------
\section{Data and Infrastructure Prerequisites}
%% ----------------------------------------------------------

No AI initiative succeeds without the data foundation beneath it. The
following are non-negotiable prerequisites before machine learning work
begins.

\spacer

\noindent\colorbox{primary!8}{%
\begin{minipage}{\dimexpr\textwidth-2\fboxsep}
  \vspace{6pt}
  \textbf{\color{primary}Communications Infrastructure as a Transformation Dependency.}\quad
  The group's existing LAN, WAN, voice, and remote rig data communications
  infrastructure --- spanning Lagos HQ, Ogharefe, Owa Aladinma, and Newcross
  E\&P rig sites --- is the physical foundation for every data initiative in
  this roadmap.

  \vspace{4pt}
  A degraded WAN link to Ogharefe does not just reduce productivity: it
  silences SCADA feeds and creates blind spots in operational monitoring.
  A failed VSAT link to a remote rig means drilling parameter data stops
  transmitting and decisions are made on stale information.

  \vspace{4pt}
  Communications infrastructure health is therefore a dependency of the
  transformation programme, not a separate IT concern. A connectivity
  audit across all sites must run in parallel with the data audit.
  \vspace{6pt}
\end{minipage}}

\spacer
\small
\begin{tabularx}{\columnwidth}{@{}lX@{}}
\toprule
\textbf{Prerequisite} & \textbf{What It Means} \\
\midrule
Communications and connectivity audit
  & Assess WAN, LAN, VSAT, and rig data link capacity, reliability, and
    redundancy across all sites (Lagos HQ, Ogharefe, Owa Aladinma, rig
    sites). Identify gaps that would prevent reliable data pipeline
    operation. Runs in parallel with the data audit. \\
Data audit
  & Assess quality, completeness, and accessibility of SAP, SCADA, and
    operational data across all three entities. \\
Data governance framework
  & Define ownership, quality standards, access controls, and retention
    policies across the group; NDPR compliance built in. \\
Sensor instrumentation --- OOGP
  & Install vibration, temperature, and pressure sensors on priority
    compressor and cryogenic units where not already in place. \\
SAP data connectors
  & Build or procure reliable, authenticated API connections to SAP
    financial, maintenance, procurement, and HR modules. \\
Azure IoT Hub configuration
  & Set up IoT Hub for OOGP sensor ingestion; configure time-series
    database (InfluxDB or TimescaleDB) for sensor data storage. \\
Satellite imagery subscription
  & NASRDA (Nigerian-first), Planet Labs, or Airbus Defence for
    AEPP corridor change detection. \\
\bottomrule
\end{tabularx}
\normalsize

\bigspacer

%% ----- Footer -----
{\small\color{darkgrey}
\textbf{Newcross and Pan Ocean Group of Companies}\enspace|\enspace
AI and Data Roadmap --- Version 2.0\enspace|\enspace May 2026\\
Companion document: \textit{Digital Transformation --- A Perspective Paper
(Newcross and Pan Ocean Group, May 2026)}.\\
All AI use cases are grounded in available data, existing infrastructure,
and Nigerian operational reality --- not research projects.
}

\end{document}
